\section{Introduction}
Electrochemical impedance spectroscopy (EIS) is a widely used method for label-free biosensing \cite{daniels2007label} because it characterizes interfacial properties, such as solution resistance ($R_s$), charge-transfer resistance ($R_{ct}$), and double-layer capacitance ($C_{dl}$), with minimal system perturbation \cite{lisdat2008use}. Unlike cyclic voltammetry, which applies a potential sweep that can drive faradaic processes, EIS typically applies a small-signal excitation (5--15~mV$_p$) around the rest potential. This allows functionalized surfaces to be measured repeatedly with limited perturbation \cite{daniels2007label, magar2021electrochemical}. To capture these dynamics, an instrument must cover a broad frequency range (0.01--0.1~Hz to 100~kHz) \cite{daniels2007label, lisdat2008use, brett2022electrochemical, lazanas2023electrochemical}. The upper decades resolve the solution-resistance plateau and the high-frequency side of the interfacial semicircle, while the lower decades capture slow interfacial and diffusional processes

The high cost and large footprint of benchtop wide-band EIS instruments limit their deployment in point-of-care, field, and educational settings \cite{pruna2018low, burgos2022tbistat, bautista2025helpstat}. While many low-cost, open-source potentiostats have been developed to address this limitation \cite{bill2023electrochemical, rowe2011cheapstat, anshori2025ganestat}, most are restricted to voltammetric methods. However, modern integrated analog front-ends, particularly the Analog Devices AD5941, now make wide-band EIS feasible on compact, low-power boards \cite{bautista2025helpstat, vamos2025hunstat2}.

Despite hardware advances, demonstrating quantitative agreement using predefined acceptance criteria remains inconsistent. As summarized in Table~\ref{tab:comparison}, validation in low-cost EIS systems relies primarily on spectral overlays, error bounds, regression, or hypothesis testing. However, these methods fall short. Significance testing ignores practical importance, accuracy bounds do not verify actual compliance, and visual overlays obscure systematic errors. Furthermore, studies rarely apply explicit criteria to flag invalid data. This work therefore focuses on distinguishing statistically significant differences from practically acceptable differences using a predefined equivalence criterion.

This work presents STEI-ZStat, an AD5941-based potentiostat costing US\$25.36 (0.2~Hz--200~kHz) on a 63.1 $\times$ 45.0~mm board. We also introduce a four-stage validation procedure: reactance-sign screening, resistive plateau checks, complex nonlinear least-squares (CNLS) parameter recovery, and equivalence testing against a commercial reference. Performance was evaluated using a precision Randles dummy cell, comparing measurements against a factory-calibrated PalmSens4 reference and the spectrum predicted from nominal component values. Board-to-board reproducibility was assessed in supplementary measurements on three assembled boards. Validation using physical biosensors is deferred to future work.


\begin{table*}[!t]
\renewcommand{\arraystretch}{1.3}
\caption{Comparison Table of Low-Cost EIS-Capable Potentiostats}
\label{tab:comparison}
\centering
\resizebox{\textwidth}{!}{
\begin{threeparttable}
\begin{tabular}{lcccc >{\raggedright\arraybackslash}p{5cm}}
\hline
\multirow{2}{*}{\textbf{Name}}
& \textbf{Freq. Range}
& \textbf{Reported Cost\tnote{a}}
& \textbf{Dimensions}
& \multirow{2}{*}{\textbf{Electroanalytical Methods}}
& \multirow{2}{*}{\textbf{Validation Methods}} \\
& \textbf{(Hz)} & \textbf{(USD)} & \textbf{(mm)} & & \\
\hline
Pruna et al.\ \cite{pruna2018low} & 0.1--10,000 & 300\tnote{b} & 120 $\times$ 70 & EIS & Regression against VMP3, $R^2= 0.98$ \\
ABE-STAT \cite{jenkins2019abe} & 0.1--100,000 & 105\tnote{c} & 89 $\times$ 74 & CV, DPV, OCP, EIS & RMS noise; potentiometric regression \\
Matsubara \cite{matsubara2021small} & 0.1--1,000,000 & 410\tnote{c,f} & Not stated & CV, SWV, CA, CP, EIS & Error bounds only \\
TBISTAT \cite{burgos2022tbistat} & 1--10,000 & 77.4\tnote{d} & Not stated & EIS & t-test; Levene; normality checks \\
ACEStat \cite{brown2022acestat} & 1--120,000 & 60\tnote{d} & 68 $\times$ 26 & CV, LSV, SWV, cSWV, CA, EIS & Randles fitting; margin of error \\
HELPStat \cite{bautista2025helpstat} & 0.15--200,000 & 50.5\tnote{e} & 61.7 $\times$ 44.4 & EIS & Paired t-test; Theil--Sen \\
HunStat2 \cite{vamos2025hunstat2} & 0.2--200,000 & 30\tnote{e,g} & 80 $\times$ 50\tnote{h} & CV, OCP, EIS & Visual overlay only \\
\textbf{STEI-ZStat (this work)} & 0.2--200,000 & 25.36\tnote{e} & 63.1 $\times$ 45.0 & EIS\tnote{i} & Reactance-sign screening; plateau checks; CNLS recovery; equivalence testing \\
\hline
\end{tabular}
\begin{tablenotes}
\item[] $^\text{a}$ Costs in USD are as reported by each source and have not been renormalized; this work's figure is per unit at a build quantity of five (May 2026). 
$^\text{b}$ Details not stated. 
$^\text{c}$ Components only. 
$^\text{d}$ Components and board fabrication. 
$^\text{e}$ Components, fabrication, and assembly.
$^\text{f}$ Calculations included a \$279 oscilloscope in the BOM.
$^\text{g}$ Recalculated from the published bill of materials and stated manufacturing costs; approximately \$27.5.
$^\text{h}$ Taken from the open-sourced board design files.
$^\text{i}$ CV, OCP, CA, SWV, and DPV are also implemented in firmware but are not experimentally validated in this work.
\end{tablenotes}
\end{threeparttable}}
\end{table*}



\begin{figure}[t]
\centering
\includegraphics[width=0.95\columnwidth]{fig1_STEI-ZStat.png}
\caption{(a) Block diagram and (b) photograph of STEI-ZStat with Randles dummy cell connected.}
\label{fig:system_design}
\end{figure}
